\documentclass{article}
\usepackage{graphicx} % Required for inserting images
\usepackage{caption} % Required for \captionof
\usepackage{amsmath}
\usepackage{float}
\usepackage{comment}
\usepackage[T1]{fontenc}
\usepackage[utf8]{inputenc}
\usepackage{lmodern}
\usepackage{microtype}

\title{Hőtani mérés}
\author{Kern Luca, Csongor Vincze}
\date{2026 Április 1}

\begin{document}

\maketitle

\section*{4. Feladat: A termosz hőkapacitásának becslése}

%hibabecslés: nem a tömegmérés relatív hibája, hanem a 2-vel való osztásnak. Annak van egy 10-15-20%-os hibája. Ezt majd át kell gondolni!

Ebben a feladatban a vizsgált termoszok belső hőkapacitását becsültük meg geometriai alapon. Ahogy azt korábban már említettük, az a hőmennyiség, amit a mérőedény fala a mérés során felvesz, alapvetően befolyásolja az összes vizsgált folyamat energiamérlegét. Ehhez megmértük az egyes termoszok teljes tömegét. Mivel a felépítésük kétrétegű, feltételeztük, hogy a belső és a külső acélfaluk ugyanolyan vastag, és mivel közöttük kiváló hőszigetelő tulajdonságú vákuum található, nincsen számottevő hőcsere a két falrész között. Emiatt a termikus folyamatokban ténylegesen részt vevő hőtároló tömeget a termosz teljes tömegének felével közelítettük. Ezt követően az edények hőkapacitását a $C = c_{\text{acél}} \cdot \frac{m_{\text{termosz}}}{2}$ összefüggés alapján számítottuk ki, a rozsdamentes acél irodalmi fajhőjének ($c_{\text{rozsdamentes acél}} = 500 \text{ J/(kg}^{\circ}\text{C)}$) felhasználásával.

\begin{table}[H]
    \centering
    \begin{tabular}{|l|c|c|}
        \hline
        Termosz & Teljes mért tömeg [g] & Becsült hőkapacitás [J/$^{\circ}$C] \\\hline
        Zöld & 430 & 107,5 \\ \hline
        Fekete & 444 & 111 \\ \hline
        Piros & 380 & 95 \\ \hline
    \end{tabular}
    \caption{Termoszok hőkapacitásának becslése}
    \label{tab:hotan_4}
\end{table}

\begin{comment}
$m_{\text{fekete}}=444$ g\\
$m_{\text{zöld}}=430$ g\\
$m_{\text{piros}}=380$ g\\
$c_{\text{rozsdamentes acél}}=500 \text{ J/(kg}^{\circ}\text{C)}$\\
Becsült hőkapacitás:\\
$C_{\text{fekete}}=111$ J/°C\\
$C_{\text{zöld}}=107,5$ J/°C\\
$C_{\text{piros}}=95$ J/°C\\
\end{comment}

\section*{5. Feladat: A víz fajhőjének mérése elektromos fűtéssel}

%fontos: előtte és utána is mérni 1-2 percig!!!
Ebben a feladatban a víz fajhőjét határoztuk meg elektromos fűtéssel. A zöld termoszunkat megtöltöttük 0,77 liter hideg vízzel, majd a folyadék folyamatos belső hőleadása (elektromos fűtés) melletti melegedését vizsgáltuk. A mérés kezdetén feljegyeztük a termoszban lévő még fűtetlen víz tömegét és kezdőhőmérsékletét. Ahhoz, hogy a fűtőtestet a méréshez javasolt stabil 50 W-os hőteljesítményen tudjuk üzemeltetni, ismernünk kellett az ellenállását. A digitális multimétert ellenállásmérő módba kapcsolva az fűtőszál ellenállása $1,23\ \Omega$-nak adódott. A minél nagyobb biztonság, valamint a garantáltan stabil teljesítményleadás érdekében a fűtőtestet egy áramgenerátorra kötöttük, és finoman beállítottuk rajta a szükséges feszültséget. A leolvasott áramerősség végül 6,46 A, az ehhez tartozó feszültség pedig 7,95 V lett (melyből pontosan számítható a villamos teljesítmény). Ezen beállítások mellett pontosan 10 percig folyamatosan fűtöttük a vizet, közben mágneses keverőt alkalmazva a homogén hőmérséklet-eloszlás elérésére.

A rögzített hőmérsékletadatokat ábrázoltuk az eltelt idő függvényében. Az így kapott monoton növő szakaszra egyenest illesztettünk (ez volt az ahol bekapcsoltuk a fűtést), amiből meghatároztuk a hőfoknövekedés meredekségét. Végül a termodinamikai energiamérleg felírásával megállapítottuk a víz fajhőjét, ehhez a termosz korábban becsült hőkapacitását is felhasználtuk. A számolás eredményeként a víz fajhőjére $4533 \text{ J/(kg}^{\circ}\text{C)}$ adódott. Ami megint nem tökéletes, de nincs olyan távol az irodalmi $4178 \text{ J/(kg}^{\circ}\text{C)}$ értéktől. A relatív hiba 8,5 $\%$, ami egy ilyen mérési beállításnál elfogadható. 

\begin{comment}
$V_{\text{víz}}=$0,77 L\\
$m_{\text{víz}}=$774 g\\
$R_{\text{fűtőtest}}=1,23 \Omega$\\
$U_{\text{fűtőtest}}=$7,95 V\\
$I_{\text{fűtőtest}}=6,46$ A\\
$T_{kezdeti}=26,4$ °C\\
$T_{veg}=36,4$ °C\\
termosz hőkapacitása az előző feladatból\\
$c_{\text{víz}}=4092 \text{ J/(kg}^{\circ}\text{C)}$\\

Ezek alapján a fajhő $c_{\text{víz}}=4532,52 \text{ J/(kg}^{\circ}\text{C)}$
-> Mi ez az elentmondás?????
\end{comment}

\begin{figure}[H]
    \centering
    \includegraphics[width=0.8\textwidth]{4_1_graf.pdf}
    \caption*{5. feladat: Fűtött víz hőmérsékletfüggése}
    \label{fig:5_1}
\end{figure}

Hibaszámítás: ..........

\section*{6. Feladat: Termosz hőkapacitásának mérése keveréssel}

Ebben a feladatban a melegvizes termosz hőkapacitását határoztuk meg kalorimetriás keveredési módszerrel (különböző hőmérsékletű víztömegeket megkeverve). A megfelelő termoszokba beleöntöttük a meleg, illetve a hideg vizet, majd precízen megmértük a tömegüket. A hidegvizes edényben 464 g, a melegvizesben pedig 467 g víz volt. Miután mindkét termoszban beállt a belső termikus egyensúly, leolvastuk a kezdeti hőmérsékleteket. Ezután -- a mágneses keverő folyamatos működtetése mellett -- óvatosan beleöntöttük a hideg vizet a melegvizes termoszba, mialatt a keverék hőmérsékletének változását is folyamatosan rögzítettük. Az így kapott mérési adatsorból grafikonokat készítettünk, meghatározva a termikus egyensúlyi értékeket. A termosz hőkapacitásának meghatározásához az alábbi összefüggést írtuk fel: \[
m_{\text{meleg víz}} \cdot c_{\text{v}} \cdot (T'_{\text{meleg víz}} - T'_{\text{közös}}) + C_{\text{termosz}} \cdot (T'_{\text{meleg víz}} - T'_{\text{közös}}) = m_{\text{hideg víz}} \cdot c_{\text{v}} \cdot (T'_{\text{közös}} - T'_{\text{hideg víz}}).
\]
A számítások elvégzéséhez feltételeztük, hogy a keverési folyamat ideje alatt a rendszer hőszigetelt, teljesen zártnak tekinthető (vagyis a környezetbe hőáramlások elhanyagolhatók), miközben a víz fajhőjének (szobahőmérsékleti) 4178 J/(kg$^{\circ}$C) irodalmi értékét vettük alapul. Ezen feltételezésekkel a zöld termosz vizsgált hőkapacitása $197,4 \text{ J/}^{\circ}\text{C}$-ra adódott. Ami ugyan több mint 2-szer akkora mint a 4. feladatban becsült érték, de a mérés pontatlanságához képest elfogadható. Itt nagyon nagy szerepet játszott az, hogy a mérést hosszú ideig végeztük, és így a rendszer nem hőszigetelt.

\begin{comment}
Ezek az első alkalom eredményei.\\
$m_{hideg}=464$ g\\
$m_{meleg}=467$ g\\
$T_{hideg}=110,4 \Omega$\\
$T_{meleg}=141,0 \Omega$\\
$T_{\text{egyensúlyi}}=125,0 \Omega$\\
$c_{\text{víz}}= 4178 \text{ J/(kg}^{\circ}\text{C)}$ \\  %irodalmi adat
$c_{\text{rozsdamentes acél}}=500 \text{ J/(kg}^{\circ}\text{C)}$\\
Zöld termosz hőkapacitása keveréssel: $C_{termosz}=197.4$\\
\end{comment}

Sajnos mivel a második laboratóriumi alkalommal a szokványos zöld helyett mi egy piros termoszt kaptunk (ami térfogatban is különbözött az eddigiektől), a kalorimetriás kalibrálást ismét meg kellett ejtenünk ezen a piros eszközön is.

A megismételt második alkalommal is teljesen ugyanezzel a metodikával dolgoztunk a piros termosznál, ekkor viszont 294 g hideg és 406 g meleg vizet öntöttünk egybe. A hideg víz kezdő hőmérséklete 22,3 $^{\circ}$C, míg a meleg vízé 85,0 $^{\circ}$C volt. A keveredés utáni egyensúly, 58,4 $^{\circ}$C értéken állt be. Viszont a fenti mérlegegyenlet visszaszámolásakor a piros termosz hőkapacitására csupán egy $23,6 \text{ J/}^{\circ}\text{C}$-os érték adódott, amely a valóságosnál vélhetőleg jóval kisebb, de legalább nem negatív. A mérést nem ismételtük meg harmadszor, a laborvezető instrukciója szerint.
\begin{comment}
Ezek a második alkalom eredményei.\\
$m_{hideg}=294$ g\\
$m_{meleg}=406$ g\\
$R_{hideg}=110,29 \Omega$\\
$R_{meleg}=134,45 \Omega$\\
$T_{hideg}=22,3$ °C\\
$T_{meleg}=85,0$ °C\\
$R_{\text{egyensúlyi}}=125,0 \Omega$\\
$T_{\text{egyensúlyi}}=58,4$ °C\\
$c_{\text{víz}}= 4178 \text{ J/(kg}^{\circ}\text{C)}$ \\  %irodalmi adat
$c_{\text{rozsdamentes acél}}=500 \text{ J/(kg}^{\circ}\text{C)}$\\
Piros termosz hőkapacitása keveréssel: $C_{termosz}=23,6$ J/K\\
\end{comment}
\begin{figure}[H]
    \centering
    \includegraphics[width=0.8\textwidth]{6_1_graf.pdf}
    \caption*{6. feladat: Zöld termosz hőkapacitásának mérése keveréssel}
    \label{fig:6_1}
\end{figure}

\begin{figure}[H]
    \centering
    \includegraphics[width=0.8\textwidth]{6_2_graf.pdf}
    \caption*{6. feladat: Piros termosz hőkapacitásának mérése keveréssel}
    \label{fig:6_1}
\end{figure}

Hibaszámítás: .........





\end{document}